\documentclass[12pt]{article}
\usepackage{float}
\usepackage[a4paper]{geometry}
\geometry{top=2.54cm, bottom=2.54cm, left=2.54cm, right=2.54cm}
\usepackage{graphicx,subfig} %LaTeX package to import graphics
\graphicspath{{images/}} %configuring the graphicx package
\usepackage{amsmath}
\usepackage{amsfonts}

\begin{document}
\section*{Oscilador armonico cuántico:}

El Hamiltoniano del oscilador armónico en 1D es:$$\left( -\frac{\hbar^2}{2m} \frac{d^2}{dx^2} + \frac{1}{2}m\omega^2 x^2 \right) \psi(x) = E \psi(x)$$Reordenando para dejar la derivada de mayor orden limpia:$$\frac{d^2\psi(x)}{dx^2} - \frac{m^2\omega^2}{\hbar^2} x^2 \psi(x) + \frac{2mE}{\hbar^2} \psi(x) = 0$$Definimos las constantes físicas estándar con dimensiones apropiadas:$a^2 = \frac{m\omega}{\hbar} \implies a^4 = \frac{m^2\omega^2}{\hbar^2}$$k^2 = \frac{2mE}{\hbar^2}$La ecuación toma la forma:$$\mathbf{\frac{d^2\psi(x)}{dx^2} - a^4 x^2 \psi(x) + k^2 \psi(x) = 0}$$Paso 2: Cambio de Variable a $y = a^2 x^2$ (Hacia la Ecuación Confluente)Al realizar la transformación de variable $y = a^2 x^2$, el operador de segunda derivada se transforma rigurosamente en:$$\frac{d^2}{dx^2} \longrightarrow 4a^2 y \frac{d^2}{dy^2} + 2a^2 \frac{d}{dy}$$Sustituyendo $\frac{d^2\psi}{dx^2}$ y $x^2 = \frac{y}{a^2}$ en la ecuación diferencial:$$\left( 4a^2 y \frac{d^2\psi}{dy^2} + 2a^2 \frac{d\psi}{dy} \right) - a^4 \left(\frac{y}{a^2}\right)\psi + k^2 \psi = 0$$Dividiendo toda la ecuación entre $4a^2$:$$y \frac{d^2\psi}{dy^2} + \frac{1}{2} \frac{d\psi}{dy} + \left( \frac{k^2}{4a^2} - \frac{y}{4} \right) \psi(y) = 0$$Evolución del coeficiente de la Energía:Definimos el parámetro adimensional de energía $\mathcal{E}$:$$\mathcal{E} \equiv \frac{k^2}{2a^2} = \frac{\frac{2mE}{\hbar^2}}{2\frac{m\omega}{\hbar}} = \mathbf{\frac{E}{\hbar\omega}}$$Sustituyendo $\frac{k^2}{4a^2} = \frac{\mathcal{E}}{2}$, obtenemos la ecuación intermedia:$$\mathbf{y \frac{d^2\psi(y)}{dy^2} + \frac{1}{2} \frac{d\psi(y)}{dy} + \left( \frac{\mathcal{E}}{2} - \frac{y}{4} \right) \psi(y) = 0}$$Paso 3: Extracción del Comportamiento Asintótico $e^{-y/2}$ (Ecuación de Kummer)Para grandes valores de $y \to \infty$, la ecuación dominante es $\psi'' - \frac{1}{4}\psi \approx 0$, sugiriendo una solución de la forma $\psi(y) = e^{-y/2} w(y)$.Sustituyendo $\psi(y) = e^{-y/2} w(y)$ y sus derivadas:$\psi' = e^{-y/2} \left( w' - \frac{1}{2}w \right)$$\psi'' = e^{-y/2} \left( w'' - w' + \frac{1}{4}w \right)$Sustituyendo estas derivadas en la ecuación anterior y cancelando el factor común $e^{-y/2}$:$$y \left( w'' - w' + \frac{1}{4}w \right) + \frac{1}{2} \left( w' - \frac{1}{2}w \right) + \left( \frac{\mathcal{E}}{2} - \frac{y}{4} \right) w = 0$$Desarrollando los términos:$$y w'' - y w' + \frac{y}{4}w + \frac{1}{2} w' - \frac{1}{4}w + \frac{\mathcal{E}}{2}w - \frac{y}{4}w = 0$$Cancelando los términos $\frac{y}{4}w$ y $-\frac{y}{4}w$:$$y w'' + \left(\frac{1}{2} - y\right) w' - \left( \frac{1}{4} - \frac{\mathcal{E}}{2} \right) w = 0$$Llegamos exactamente a la Ecuación Hipergeométrica Confluente (de Kummer):$$\mathbf{y \frac{d^2 w(y)}{dy^2} + (b - y) \frac{dw(y)}{dy} - a \, w(y) = 0}$$donde los parámetros corresponden a:$b = \frac{1}{2}$$a = \frac{1}{4} - \frac{\mathcal{E}}{2} = \mathbf{\frac{1}{4} - \frac{E}{2\hbar\omega}}$Paso 4: Transformación a la Variable Espacial $\xi = ax = \sqrt{y}$ (Ecuación de Hermite)Si en lugar de $y$ usamos la coordenada espacial adimensional $\xi = ax = \sqrt{\frac{m\omega}{\hbar}} x$, se extrae la envolvente gaussiana $\psi(\xi) = e^{-\xi^2/2} H(\xi)$.Calculando las derivadas de $\psi(\xi) = e^{-\xi^2/2} H(\xi)$ respecto a $\xi$:$\frac{d\psi}{d\xi} = e^{-\xi^2/2} \left( \frac{dH}{d\xi} - \xi H \right)$$\frac{d^2\psi}{d\xi^2} = e^{-\xi^2/2} \left( \frac{d^2H}{d\xi^2} - 2\xi \frac{dH}{d\xi} + (\xi^2 - 1) H \right)$Expresando la ecuación de Schrödinger en la variable $\xi$ ($x^2 = \frac{\xi^2}{a^2}$ y $k^2/a^2 = 2\mathcal{E}$):$$\frac{d^2\psi}{d\xi^2} - \xi^2 \psi + 2\mathcal{E} \psi = 0$$Sustituyendo $\psi(\xi)$ y sus derivadas, y cancelando $e^{-\xi^2/2}$:$$\left( \frac{d^2H}{d\xi^2} - 2\xi \frac{dH}{d\xi} + (\xi^2 - 1) H \right) - \xi^2 H + 2\mathcal{E} H = 0$$Cancelando $\xi^2 H$ y $-\xi^2 H$:$$\mathbf{\frac{d^2 H(\xi)}{d\xi^2} - 2\xi \frac{dH(\xi)}{d\xi} + (2\mathcal{E} - 1) H(\xi) = 0}$$

\end{document}
